\section{Methodology: Slot-Kinetics}
\label{sec:method}

We introduce \textbf{Tenant-Decoupled Stateful Routing (TDSR)} as a robust and highly efficient overarching framework for stateful test-time model ensembling under interleaved, multi-tenant query streams, powered by our underlying \textbf{Slot-Kinetics} state decoupling mechanism. We first formalize feature extraction and subspace coordinate projection, detail the multi-tenant state decoupling architecture, formulate the state updates, and present a systems-level complexity and deployment analysis.

\subsection{Subspace Coordinate Projection}
Consider a deep neural network of $L$ layers. The boundary layer $l_{\text{route}} < L$ is frozen and designated for routing feature extraction. At inference step $t$, the input query passes through the early boundary layers, yielding intermediate activation vector $h_t^{(l_{\text{route}})} \in \mathbb{R}^D$, where $D$ is the hidden dimension. To extract task coordinates, we first unit-normalize the activation vector to ensure scale invariance:
\begin{equation}
\tilde{h}_t^{(l_{\text{route}})} = \frac{h_t^{(l_{\text{route}})}}{\|h_t^{(l_{\text{route}})}\|_2 + \epsilon_{\text{norm}}}
\end{equation}
where $\epsilon_{\text{norm}} = 10^{-6}$ is a numerical stabilizer. 

For a fleet of $K$ experts, we pre-compute task-specific projection matrices $V_{k, d} \in \mathbb{R}^{D \times d}$ offline via Principal Component Analysis (PCA) on task-specific calibration activations. At test-time, the normalized activation $\tilde{h}_t^{(l_{\text{route}})}$ is projected onto these task-specific subspaces. The coordinate vector $\mathbf{e}_t = [e_{0, t}, \dots, e_{K-1, t}]^T \in [0, 1]^K$ is computed as the $L_2$-norm of the projected activations:
\begin{equation}
e_{k, t} = \|V_{k, d}^T \tilde{h}_t^{(l_{\text{route}})}\|_2 \quad \forall k \in \{0, \dots, K-1\}
\end{equation}
The vector $\mathbf{e}_t$ represents the current input's affinity to each task-specific expert.

\subsection{Multi-Tenant State Decoupling}
Existing stateful routers (e.g., PAC-Kinetics \cite{pac_zca_2026}) maintain a single global routing state $\mathbf{s}_t \in \mathbb{R}^K$. Under interleaved workloads, this causes cross-tenant state contamination. To resolve this, TDSR maintains a pool of $M$ active \emph{state slots} $\mathcal{S} = \{\mathbf{s}_0, \dots, \mathbf{s}_{M-1}\}$ and $M$ slot centroids $\mathcal{C} = \{\mathbf{c}_0, \dots, \mathbf{c}_{M-1}\}$, where $\mathbf{s}_m \in \mathbb{R}^K$ and $\mathbf{c}_m \in [0, 1]^K$ for $m \in \{0, \dots, M-1\}$. At step $t$, the current query must be routed to its respective virtual slot index $m^*_t \in \{0,\dots, M-1\}$. We propose two pragmatic session routing modes:

\paragraph{Scenario A: Explicit Session Tagging.}
If the serving framework (e.g., S-LoRA \cite{s-lora} or Punica \cite{chen2023punica}) is metadata-aware, the central scheduler supplies a session or tenant ID $u_t \in \{0, \dots, M-1\}$ alongside the input batch. Under this mode, the winning slot index is set directly:
\begin{equation}
m^*_t = u_t
\end{equation}
This achieves perfect session isolation with absolutely zero computational overhead.

\paragraph{Implicit Tagless Clustering.}
In lightweight or edge deployments where user metadata is unavailable, the router must dynamically infer the session context on-the-fly. Since the coordinate vector $\mathbf{e}_t$ represents the query's activation-space projection onto $K$ offline-defined expert subspaces, the coordinate axes themselves represent the ideal activation profiles. To achieve robust, non-drifting slot assignment, we initialize and fix the slot centroids $\mathcal{C}$ as orthogonal coordinate detector vectors ($\mathbf{c}_{m, m} = 1.0$). This completely eliminates centroid drift and runaway slot attraction (clustering collapse), ensuring perfect unsupervised slot specialization. 

At each step, we compute the cosine similarity between the coordinate vector $\mathbf{e}_t$ and each fixed task detector centroid $\mathbf{c}_m$:
\begin{equation}
\label{eq:cosine_sim}
\text{Sim}(\mathbf{e}_t, \mathbf{c}_m) = \frac{\mathbf{e}_t^T \mathbf{c}_m}{\|\mathbf{e}_t\|_2 \|\mathbf{c}_m\|_2 + \epsilon_{\text{sim}}}
\end{equation}
where $\epsilon_{\text{sim}} = 10^{-6}$. The query is then routed to the slot with the highest similarity:
\begin{equation}
\label{eq:implicit_routing}
m^*_t = \arg\max_{m \in \{0, \dots, M-1\}} \text{Sim}(\mathbf{e}_t, \mathbf{c}_m)
\end{equation}
To provide crucial scientific and systems transparency, we highlight that because the slot centroids $\mathcal{C}$ are initialized and fixed as standard orthogonal coordinate detector vectors (i.e., $\mathbf{c}_m$ is the $m$-th standard basis vector where $\mathbf{c}_{m, m} = 1.0$ and $\mathbf{c}_{m, j} = 0.0$ for $j \neq m$), the online cosine similarity in Eq.~\ref{eq:cosine_sim} simplifies mathematically to:
\begin{equation}
\label{eq:cosine_sim_simple}
\text{Sim}(\mathbf{e}_t, \mathbf{c}_m) = \frac{e_{m, t}}{\|\mathbf{e}_t\|_2 \cdot 1.0 + \epsilon_{\text{sim}}}
\end{equation}
Because the denominator $\|\mathbf{e}_t\|_2 + \epsilon_{\text{sim}}$ is a positive scalar constant identical for all candidate slots $m \in \{0, \dots, M-1\}$ at any given step $t$, the argmax similarity operation in Eq.~\ref{eq:implicit_routing} simplifies exactly to:
\begin{equation}
\label{eq:argmax_coord}
m^*_t = \arg\max_{m \in \{0, \dots, M-1\}} e_{m, t}
\end{equation}
Thus, while the formulation is conceptually grounded in online cosine similarity clustering against a fixed basis, the actual system-level execution reduces to a highly efficient and simple \textbf{coordinate-argmax assignment} (selecting the slot corresponding to the task with the maximum activation coordinate). This mathematical simplification is of immense practical significance for high-throughput serving: it completely bypasses vector dot-products, norm computations, and divisions, allowing slot assignment to be implemented as a single, sub-nanosecond integer index lookup over the coordinates.

We emphasize that keeping the centroids fixed as coordinate detectors does not constitute a label-leaking Oracle. Since the task coordinates are computed entirely from model-side activations projected onto offline PCA expert spaces, the assignment process is completely unsupervised and requires zero ground-truth label feedback during deployment.

\paragraph{The Slot-Tenant-Task Triad: Virtual Task Caching.}
Under the implicit tagless mode, because centroids are fixed as orthogonal task detectors, the routing mapping in Eq~\ref{eq:implicit_routing} groups incoming queries based on their *task affinity* rather than physical *tenant identity*. In this configuration, the virtual slots specialize in specific task domains, transforming the system into a \textbf{Task-Decoupled Stateful Router}. 
From a systems-engineering perspective, this represents an exceptionally elegant design pattern. Instead of allocating physical state slots to every individual user session (which scales as $\mathcal{O}(M)$ and introduces massive memory pools for high-concurrency cloud servers), the router maps requests to a small, fixed pool of $K$ task-specific virtual slots (where $K \ll M$). This completely decouples the slot memory from the number of active users, allowing a shared cloud gateway to scale to an arbitrary number of concurrent tenants while still completely preventing state contamination. Since queries belonging to the same task from different users are routed to the same specialized slot, they share a mutually beneficial, smoothed task-specific context rather than experiencing cross-tenant interference from unrelated tasks.

\subsection{Stateful Recurrence and Tenant-Specific Session-Step Decay}
Once the winning slot index $m^*_t$ is determined, we update the state pool $\mathcal{S}$ as follows:
\begin{itemize}
    \item \textbf{Active Slot Update (First-Order Recurrence):} The active slot's concentration state undergoes a recurrent update:
    \begin{equation}
    \mathbf{s}_{m^*_t, t} = \mathbf{A} \mathbf{s}_{m^*_t, t-1} + W \mathbf{e}_t
    \end{equation}
    where $\mathbf{A} = \text{diag}(a_0, \dots, a_{K-1})$ with $a_k = \sigma(u_k) \in (0, 1)$ is the task-specific state retention matrix, and $W \in \mathbb{R}^{K \times K}$ is the unconstrained coordinate injection coupling matrix.
    \item \textbf{Inactive Slots (Logical Session-Step Decay):} While global-step exponential decay $\mathbf{s}_{m, t} = A_{\text{decay}} \mathbf{s}_{m, t-1}$ is simple, it poses a scalability boundary in high-throughput cloud environments: if the server hosts $M=1000$ active tenants, a specific tenant's queries will be highly sparse in global steps, causing their stateful memory to be prematurely washed out by other tenants' traffic.
    
    To resolve this, we formulate \textbf{Tenant-Specific Session-Step Decay} (or logical step decay). Under this policy, inactive slots do not decay on global steps. Instead, a slot's state is only decayed on its own logical session steps (i.e., on its own local clock when active) or when a physical temporal timeout is exceeded. Mathematically:
    \begin{equation}
    \mathbf{s}_{m, t} = \begin{cases} \mathbf{A} \mathbf{s}_{m, t-1} + W \mathbf{e}_t & \text{if } m = m^*_t \\ A_{\text{decay}}^{\Delta t_m} \mathbf{s}_{m, t-1} & \text{if } m \ne m^*_t \end{cases}
    \end{equation}
    where $\Delta t_m \ge 0$ is the elapsed logical steps of tenant $m$. In our interleaved evaluation, we set $\Delta t_m = 0$ for inactive slots (i.e., they hold their state perfectly between active queries), completely resolving the state washout bottleneck regardless of the interleaving scale $M$.
    
    To reconcile the trade-off between sequence isolation during active serving and memory reclamation of obsolete sessions, we deploy a \textbf{Dual-Clock Decay} policy (fully detailed in Section 3.6). Under this policy, inactive slots use a logical clock during active interleaved query streams (setting $\Delta t_m = 0$ to prevent state washout), while a secondary physical wall-clock timer runs in the background. If a slot remains inactive for longer than a specified physical timeout (e.g., 5 seconds), its state is exponentially decayed to zero or evicted. This dual-clock formulation resolves the logical contradiction of holding state indefinitely, providing a mathematically sound and self-cleaning memory footprint.
\end{itemize}

\subsection{Gibbs Softmax Policy Mapping}
The ensembling weights $\boldsymbol{\alpha}_t \in \Delta^{K-1}$ are computed using the winning slot's updated state $\mathbf{s}_{m^*_t, t}$ via the multi-temperature Gibbs Softmax policy:
\begin{equation}
\alpha_{k, t} = \frac{\exp(s_{m^*_t, k, t}/\tau_k)}{\sum_{j=1}^K \exp(s_{m^*_t, j, t}/\tau_j)}
\end{equation}
where the task-specific temperatures are defined as $\tau_k = e^{w_k} + \tau_{\min}$ with learnable log-temperatures $\mathbf{w}$ and a strict minimum temperature threshold $\tau_{\min} = 0.01$ to guarantee Lipschitz continuity.

Finally, for all dynamic blending layers $l \in [l_{\text{route}} + 1, L]$, intermediate representations are dynamically blended:
\begin{equation}
h_t^{(l)} = h_t^{(l-1)} + \sum_{k=1}^K \alpha_{k, t} \gamma_V (v'_k - h_t^{(l-1)})
\end{equation}
where $v'_k$ is the output of expert $k$ at layer $l$ and $\gamma_V = 0.05$ is the expert scaling coefficient. Downstream classification is evaluated at Layer $L$ via negative squared Euclidean distance:
\begin{equation}
\text{logits}_{t, k} = - \| h_t^{(L)} - v'_k \|_2^2 + b_k
\end{equation}
where $b_k$ is the task bias.

\subsection{The Step-by-Step TDSR Pipeline}
To provide a complete, systems-level view of the framework, we detail the lifecycle of a single incoming query $h_t^{(l_{\text{route}})}$ in Algorithm \ref{alg:tdsr}.

\begin{algorithm}[tb]
\caption{Tenant-Decoupled Stateful Routing (TDSR)}
\label{alg:tdsr}
\begin{algorithmic}[1]
\STATE {\bfseries Input:} Incoming query activation $h_t^{(l_{\text{route}})}$ at routing layer $l_{\text{route}}$, tenant metadata index $u_t$ (for explicit mode), active states $\mathcal{S} = \{\mathbf{s}_0, \dots, \mathbf{s}_{M-1}\}$, slot centroids $\mathcal{C} = \{\mathbf{c}_0, \dots, \mathbf{c}_{M-1}\}$, learnable parameters $\mathbf{A}$, $W$, log-temperatures $\mathbf{w}$.
\STATE {\bfseries Output:} Ensembling weights $\boldsymbol{\alpha}_t$, next layer representation $h_t^{(l_{\text{route}} + 1)}$
\STATE \COMMENT{// Step 1: Normalize and Project Activations}
\STATE $\tilde{h}_t^{(l_{\text{route}})} \leftarrow \frac{h_t^{(l_{\text{route}})}}{\|h_t^{(l_{\text{route}})}\|_2 + \epsilon_{\text{norm}}}$
\FOR{$k=0$ {\bfseries to} $K-1$}
\STATE $e_{k, t} \leftarrow \|V_{k, d}^T \tilde{h}_t^{(l_{\text{route}})}\|_2$
\ENDFOR
\STATE $\mathbf{e}_t \leftarrow [e_{0, t}, \dots, e_{K-1, t}]^T$
\STATE \COMMENT{// Step 2: Route Query to Slot}
\IF{Metadata-aware (Explicit Tagging)}
\STATE $m^*_t \leftarrow u_t$
\ELSE
\STATE $m^*_t \leftarrow \arg\max_{m \in \{0, \dots, M-1\}} \frac{\mathbf{e}_t^T \mathbf{c}_m}{\|\mathbf{e}_t\|_2 \|\mathbf{c}_m\|_2 + \epsilon_{\text{sim}}}$
\ENDIF
\STATE \COMMENT{// Step 3: Update Active and Inactive States (Decoupled Recurrence)}
\FOR{$m=0$ {\bfseries to} $M-1$}
\IF{$m = m^*_t$}
\STATE $\mathbf{s}_{m, t} \leftarrow \mathbf{A} \mathbf{s}_{m, t-1} + W \mathbf{e}_t$
\ELSE
\STATE $\mathbf{s}_{m, t} \leftarrow \mathbf{s}_{m, t-1}$ \COMMENT{// Tenant-Specific Session-Step Decay ($\Delta t_m = 0$)}
\ENDIF
\ENDFOR
\STATE \COMMENT{// Step 4: Map to Ensembling Weights}
\FOR{$k=0$ {\bfseries to} $K-1$}
\STATE $\tau_k \leftarrow e^{w_k} + \tau_{\min}$
\STATE $\alpha_{k, t} \leftarrow \frac{\exp(s_{m^*_t, k, t}/\tau_k)}{\sum_{j=1}^K \exp(s_{m^*_t, j, t}/\tau_j)}$
\ENDFOR
\STATE $\boldsymbol{\alpha}_t \leftarrow [\alpha_{0, t}, \dots, \alpha_{K-1, t}]^T$
\STATE \COMMENT{// Step 5: Dynamic Blending}
\STATE $h_t^{(l_{\text{route}} + 1)} \leftarrow h_t^{(l_{\text{route}})} + \sum_{k=0}^{K-1} \alpha_{k, t} \gamma_V (v'_k - h_t^{(l_{\text{route}})})$
\end{algorithmic}
\end{algorithm}

\subsection{Practical Systems-Level Analysis}
The Pragmatist researcher must look past mathematical formulations to evaluate actual systems-level feasibility. We highlight two key practical characteristics of our Slot-Kinetics design:

\paragraph{Sub-Microsecond Latency and Register-Level Footprint.}
In production, storing states in traditional databases (e.g., Redis) or memory pools introduces network or disk bottlenecks. Our $M \times K$ state pool represents an exceptionally small footprint. For $M=4$ tenants and $K=4$ tasks, the active states and centroids are stored as a mere $4 \times 4$ floating-point array (64 bytes). This allows the entire state pool to be pinned directly in CPU/GPU registers or fast L1 cache. Updating states and computing similarities involves basic vectorized linear operations, completing in **less than 1.5 microseconds** per query, adding virtually zero overhead compared to the millisecond-scale execution of deep Transformer layers.

\paragraph{Self-Cleaning Memory via Dual-Clock Decay.}
A common failure mode in multi-session servers is memory leak, where obsolete user sessions occupy RAM indefinitely unless cleared by complex garbage collection. To prevent this, our production-ready design implements a robust \emph{dual-clock decay} policy. During active serving blocks, the inactive slots are governed by logical session-step decay ($\Delta t_m = 0$), holding their concentration state constant to prevent state washout from interleaved queries of other tenants. However, during idle periods, a secondary physical wall-clock timer is used: if a tenant session remains inactive for longer than a specified wall-clock timeout (e.g., 5 seconds), its slot is exponentially decayed to zero: $\mathbf{s}_{m} \leftarrow A_{\text{decay}} \mathbf{s}_{m}$. This dual-clock policy provides a robust, self-cleaning session memory pool that naturally purges obsolete states and prevents memory leaks, without any garbage-collection overhead or active sequence washout.
